\documentclass[letterpaper, 11pt]{article}
\usepackage{latexsym}
\usepackage[empty]{fullpage}
\usepackage{titlesec}
\usepackage{marvosym}
\usepackage[usenames, dvipsnames]{color}
\usepackage{verbatim}
\usepackage{enumitem}
\usepackage[hidelinks]{hyperref}
\usepackage{fancyhdr}
\usepackage[english]{babel}
\usepackage{tabularx}
\usepackage{fontawesome5}
\usepackage{multicol}
\setlength{\multicolsep}{-3.0pt}
\setlength{\columnsep}{-1pt}
\input{glyphtounicode}
\usepackage{vwcol}
\usepackage{comment}
\usepackage{hyperref}
\hypersetup{
	colorlinks=true,
	linkcolor=black,
	filecolor=black,
	citecolor=black,
	urlcolor=RoyalBlue,
}

%----------FONT OPTIONS----------
% sans-serif
% \usepackage[sfdefault]{FiraSans}
% \usepackage[sfdefault]{roboto}
% \usepackage[sfdefault]{noto-sans}
% \usepackage[default]{sourcesanspro}

% serif
% \usepackage{CormorantGaramond}
% \usepackage{charter}

\pagestyle{fancy}
\fancyhf{} % clear all header and footer fields
\fancyfoot{}
\renewcommand{\headrulewidth}{0pt}
\renewcommand{\footrulewidth}{0pt}

% Adjust margins
\addtolength{\oddsidemargin}{-0.6in}
\addtolength{\evensidemargin}{-0.5in}
\addtolength{\textwidth}{1.19in}
\addtolength{\topmargin}{-.5in}
\addtolength{\textheight}{1.4in}

\urlstyle{same}

\raggedbottom
\raggedright
\setlength{\tabcolsep}{0in}

% Sections formatting
\titleformat{\section}{ \vspace{-4pt}\scshape\raggedright\large\bfseries }{}{0em}{}[
\color{black}
\titlerule
\vspace{-5pt}
]

% Ensure that generate pdf is machine readable/ATS parsable
\pdfgentounicode=1

%-------------------------
% Custom commands
\newcommand{\resumeItem}[1]{ \item\small{ {#1 \vspace{-2pt}} } }
%\newcommand{\classesList}[4]{ \item\small{ {#1 #2 #3 #4 \vspace{-2pt}} } }

\newcommand{\resumeSubheading}[4]{
\vspace{-1pt}
\item
\begin{tabular*}{1.0\textwidth}[t]{l@{\extracolsep{\fill}}r}
	\textbf{#1}       & {\small #2} \\
	{\small#3} & \textit{\small #4} \\
\end{tabular*}
\vspace{-10pt}
}

\newcommand{\resumeProjectHeading}[2]{
\vspace{-1pt}
\item
\begin{tabular*}{1.0\textwidth}[t]{l@{\extracolsep{\fill}}r}
	\textbf{#1}  & {\small #2} \\
\end{tabular*}
\vspace{-20pt}
}

% \newcommand{\resumeSubSubheading}[2]{ \item
% \begin{tabular*}{0.97\textwidth}{l@{\extracolsep{\fill}}r}
% 	\textit{\small#1} & \textit{\small #2} \\
% \end{tabular*}
% \vspace{-7pt}
% }

\newcommand{\resumeSubItem}[1]{\resumeItem{#1}
\vspace{-4pt}}

\renewcommand{\labelitemi}{$\vcenter{\hbox{\tiny$\bullet$}}$}
\renewcommand{\labelitemii}{$\vcenter{\hbox{\tiny$\bullet$}}$}

\newcommand{\resumeSubHeadingListStart}{\begin{itemize}[leftmargin=0.0in, label={}]}
\newcommand{\resumeSubHeadingListEnd}{\end{itemize}}
\newcommand{\resumeItemListStart}{\begin{itemize}[leftmargin=0.25in]}
\newcommand{\resumeItemListEnd}{\end{itemize}
\vspace{-5pt}}

%-------------------------------------------
%%%%%%  RESUME STARTS HERE  %%%%%%%%%%%%%%%%%%%%%%%%%%%%

\begin{document}
	
%----------HEADING----------
\begin{center}
    {\Huge \scshape Gawtam Chithra Ramesh} \\
    \vspace{3pt}
    % \href{https://www.kth.se/}{KTH Royal Institute of Technology} \\
    \vspace{1pt}
    \small 
    {
    \raisebox{-0.1\height}{\faPhone} +46-76541 8589 ~
    \href{gawtamcr3@gmail.com}{\raisebox{-0.2\height}{\faEnvelope}} {gawtamcr3@gmail.com} ~
    \href{https://linkedin.com/in/gawtamcr/}{\raisebox{-0.2\height}{\faLinkedin} \textcolor{black}{gawtamcr} } ~
    \href{https://gawtamcr.github.io/}{\raisebox{-0.2\height}{\faGithub} \textcolor{black}{gawtamcr}}
    }
    \vspace{-8pt}
\end{center}

%-----------EDUCATION-----------
\section{Education}
    \resumeSubHeadingListStart 
    \resumeSubheading {KTH Royal Institute of Technology}{Aug. 2024 -- Jul. 2026}{Master of Science in \href{https://www.kth.se/en/studies/master/systems-control-robotics/msc-systems-control-and-robotics-1.8733}{Systems, Controls, and Robotics} - 4.62 / 5.0}{Stockholm, Sweden}
    
    \resumeSubheading {Indian Institute of Technology Madras}{Aug. 2019 -- Jul. 2024}{Dual Degree (B.Tech in \href{https://ed.iitm.ac.in/about.html}{Engineering Design} and M.Tech in \href{https://ed.iitm.ac.in/~robotics_lab/}{Robotics}) - 7.9 / 10}{Chennai, India}

    \resumeSubHeadingListEnd

%-----------PUBLICATIONS-----------
\section{Publications}
\begin{itemize}[leftmargin=0.15in]
    \setlength{\itemsep}{2pt}
    \item \textbf{G.~Chithra~Ramesh}, et al. ``Robust Object-layer Construction of 3D Scene Graphs Using Instance Segmentation.'' \textit{International Conference on Image Processing and Vision Engineering (IMPROVE 2026)}. (Accepted, Long Oral, Nominated for Best Paper Award)
    \item \textbf{G.~Chithra~Ramesh}, et al. ``Synchronized RGB-D Inpainting for Privacy-Aware 3D Scene Graph Construction.'' \textit{International Conference on Pattern Recognition (ICPR 2026)}. (Under Review)
    \item \textbf{G.~Chithra~Ramesh}, D.~Blanco-Mulero, Y.~Dong, J.~Borr\`as Sol, C.~Torras, F.~T.~Pokorny. ``Scene Understanding in Deformable Object Manipulation via Taxonomy-Guided Vision-Language Models.'' \textit{IROS 2025 Workshop on Robotic Manipulation of Deformable Objects (ROMADO)}. (Accepted)
\end{itemize}

%-----------Scholastic and other related Achievements---------------
\section{Scholastic Achievements}
\begin{itemize}[leftmargin=0.15in]
    \setlength{\itemsep}{-0.3em}
    \item Awarded (among 30 of 800+) the IIT Madras Young Research Fellowship based on remarkable research potential.
    \item Secured 97.64 Percentile in JEE Advanced \& 99.10 Percentile in JEE Mains out of 1.5 Million across India.
    \item Secured an All India Rank of 52 out of 50k in Undergraduate Common Entrance Examination for Design.
\end{itemize}

%-----------EXPERIENCE-----------
\section{Professional Experience}
\resumeSubHeadingListStart

    \resumeSubheading{Master's Thesis on Generative Planning for Long-Horizon Manipulation}{(Ongoing) Jan 2026 - Jun 2026}{\textbf{ABB Robotics}, Mentors - 
    \href{https://www.linkedin.com/in/matthewwilliamlock/}{Matthew Lock} and \href{https://www.linkedin.com/in/jonathan-styrud-99385864/?locale=en_US}{Jonathan Styrud}}{Västerås, Sweden} 
    \resumeItemListStart 
        \resumeItem{Designing an inference-time guidance method for a pretrained flow-matching trajectory planner, allowing a single learned motion prior to satisfy arbitrary task specifications without retraining or per-task expert demonstrations.}
        \resumeItem{Using Signal Temporal Logic with STLCG++ as the specification language, enabling differentiable robustness gradients to steer generation toward compositional, temporally-extended objectives.}
        \resumeItem{Validating on Maze2D and OGBench long-horizon benchmarks; benchmarking head-to-head against state-of-the-art STL planners (DAG-STL, TeLoGraF).}
        \resumeItem{Collaborating with ABB Robotics R\&D on transfer of the approach to long-horizon 6-DoF manipulation use cases.}
    \resumeItemListEnd
    
    \resumeSubheading{Device Research Intern on 3D Scene Graph}{Jul 2025 - Dec 2025}{\textbf{Ericsson} Research, Mentor - 
    \href{https://se.linkedin.com/in/p\%C3\%BCren-g\%C3\%BCler-76457317}{{Puren Guler}} and 
    \href{https://se.linkedin.com/in/hcaltenco}{Hector Caltenco}}{Lund, Sweden} 
    \resumeItemListStart 
    \resumeItem{First-author paper accepted as Long Oral and nominated for Best Paper Award at \textbf{IMPROVE 2026}; second first-author paper under review at \textbf{ICPR 2026}.}
    \resumeItem{Developed a robust 3DSG pipeline by integrating instance segmentation and tracking into a SOTA framework (Hydra), reducing sensitivity to parameter tuning across diverse environments.}
    \resumeItem{Designed a synchronized RGB-D inpainting module for privacy-aware 3DSG construction, maintaining consistency across RGB and depth channels.}
    \resumeItem{Analyzed 3DSG frameworks to identify research gaps in object merging, parameter sensitivity, and privacy protection.}
    \resumeItem{Benchmarked SOTA inpainting algorithms to balance reconstruction fidelity with real-time performance.}
\resumeItemListEnd

    \resumeSubheading{Graduate Robotics Intern on Modular Collaborative Robots}{Dec 2022 - Jul 2023}{\textbf{Systemantics} India Pvt. Ltd., Mentor - \href{https://www.linkedin.com/in/jagannathraju/}{{Jagannath Raju}}}{Bangalore, India}
    \resumeItemListStart 
        % \resumeItem{Implemented backward-chained Behavior Trees for long-horizon 6DOF manipulator planning using ROS2 and Moveit2.}
        % \resumeItem{Implemented a control synthesis pipeline for a 6DOF manipulator using Behavior Trees to enable robust, long-horizon trajectory planning.}
        % \resumeItem{Designed and implemented a long-horizon trajectory planning system for a 6DOF manipulator using Behavior Trees,}
        \resumeItem{Developed backward-chained Behavior Trees for 6DOF manipulator to enable robust, long-horizon trajectory planning}
        \resumeItem{Designed and implemented a 6-DOF trajectory generation and control pipeline in C++, achieving real-time joint actuation via low-latency communication with motor controllers using D-Bus and SocketCAN.}
        % \resumeItem{Developed a S-Curve trajectory generator for the arm, streaming real-time joint commands directly to motor controllers via DBUS and SocketCAN.}
        % \resumeItem{Developing an disturbance observer and admittance control model for all the 6 joints of the robotic arm}
    \resumeItemListEnd
    \newpage
    \resumeSubheading{Tested Multi Object Detection and Tracking Models}{June 2021 -- July 2021}{Blurgs Research Labs Pvt. Ltd.}{Remote, India} 
    \resumeItemListStart 
        \resumeItem{Analyzed 5 State of the art Multi object tracking networks for startup's first product for The Indian Navy}
        \resumeItem{Trained and tested various MOT networks like \textbf{Siam-MOT} and \textbf{FairMOT} for Surveillance with Drones and CCTV.}
        \resumeItem{Used \textbf{FairMOT} to track and get the trail of multiple moving objects and Re-identify objects in real-time.}
    \resumeItemListEnd

    
    % \resumeProjectHeading{Graduate Teaching Assistant $|$ \textnormal{\small Field and Service Robotics (ID4060 - IITM)} } {Jan 2025 -- Feb 2025}
    \resumeSubheading {Graduate Teaching Assistant}{Aug 2023 -- Nov 2023}{Field and Service Robotics (ID4060), {Mentor - \href{https://bijosebastian.wordpress.com/}{Bijo Sebastian}}}{IIT Madras, India}
    \resumeItemListStart 
        \resumeItem{Involved in design and evaluation of assignments based on ROS and CoppeliaSim for 40 students.}
        \resumeItem{Held meetings with small groups of students with diverse backgrounds to review their assignments and projects.}
    \resumeItemListEnd
    
\resumeSubHeadingListEnd
%-----------EXPERIENCE-----------
\section{Research Experience}
\resumeSubHeadingListStart
 
    \resumeSubheading{Graduate Robotics Research on Deformable Object Manipulation}{Jan 2025 - Nov 2025}{\href{https://www.kth.se/is/rpl}{Robotics Perception and Learning}, Advisor - 
    % \href{https://www.kth.se/is/rpl}{David Blanco-Mulero} 
    %\href{https://yifeidong0.github.io/}{\emph{Yifei Dong}} and 
    \href{https://www.kth.se/is/rpl}{Florian T. Pokorny}}{KTH, Sweden} 
    \resumeItemListStart 
        % \resumeItem{Deformable object manipulation planning through taxonomy-based abstractions and foundation models}
        % \resumeItem{Developed a framework integrating pretrained-VLMs with a structured taxonomy to enhance DOM.}
        % \resumeItem{Conducted in-depth analysis of VLM performance, identifying challenges in DOM scene understanding and proposing solutions like refined taxonomy and multi-modal inputs; }
        % \resumeItem{Submitted and presented the result at ROMADO, a \textbf{IROS2025 Workshop}.}
        % \resumeItem{Designed and implemented a simulation pipeline for task execution using PyBullet, with plans for real-world validation 
        % %to assess sim-to-real transferability.
        % }

        % \resumeItem{Integrated taxonomy-guided vision-language models (Gemini, Qwen) to interpret DOM scenes and map 
        % structured outputs to 
        % motion primitives}
        % \resumeItem{Built prompt–taxonomy framework to generate structured codes and textual justifications for robot-parsable perception}
        \resumeItem{Built prompt–taxonomy framework to generate structured, robot-parsable specifications from high-level prompts, demonstrating a pipeline for translating natural-language instructions into structured outputs for downstream manipulation.}
        \resumeItem{Performed quantitative analysis to identify failure modes and guide iterative improvements to the taxonomy and prompts}
        % \resumeItem{Integrated multi-modal data and temporal cues to enhance the perception pipeline's robustness.}
        \resumeItem{Authored and submitted the results at ROMADO, an \textbf{IROS2025} workshop.}
    \resumeItemListEnd

    \resumeSubheading{Dual Degree Project on Structure from Motion}{Jan 2024 - May 2024}{\href{https://niravatgit.github.io/INSPIRELab/}{INSPIRE Lab}, Advisor - \href{https://niravatgit.github.io/patelnirav.com/index.html}{{Nirav Patel}}}{IIT Madras, India}
    \resumeItemListStart 
        \resumeItem{Developed a 3D perception system using COLMAP (Structure from Motion) from monocular RGB images to build cost-efficient human torso reconstruction for more accessible ultrasound.}
        \resumeItem{Developed an autonomous ultrasound scanning system using ROS, integrating a 5-DoF robotic arm with a single RGB camera to achieve cost-effective 3D perception for medical imaging.}
        \resumeItem{Implemented motion planning and control using ROS and MoveIt, and simulated the 5DOF manipulation in Gazebo.}
        \resumeItem{Integrated perception algorithms within ROS, processing 3D point clouds from RGB images to enable accurate trajectory planning for the ultrasound probe.}
    \resumeItemListEnd

    \resumeSubheading{Offline And Online Motion Planning \& Control of Unmanned Aerial Vehicles}{Sep 2021 -  Apr 2022}{\href{https://yrf.iitm.ac.in/}{Young Research Fellowship}, Advisor - \href{http://www.ae.iitm.ac.in/~satadal/index.html}{\emph{Prof. Satadal Ghosh}}}{IIT Madras, India}
    \resumeItemListStart
        \resumeItem{Developed a Target Driven Motion Planning and Obstacle Avoidance algorithm for UAVs.}
        \resumeItem{Used Proportional Navigation and Acceleration Velocity obstacles in 2D field to achieve 20\% faster and more computer efficient algorithm compared classical aerospace algorithms.}
    \resumeItemListEnd

\resumeSubHeadingListEnd



%-----------PROJECTS-----------
\section{Course Projects}
\resumeSubHeadingListStart

   \resumeProjectHeading{Robot Learning and Embodied AI $|$ \textnormal{\small DD2601 - KTH}}{Sep 2025 - Oct 2025}
    \resumeItemListStart 
        \resumeItem{Implemented a 3D semantic object mapping pipeline using RGB-D data from ARKitScenes, integrating OwlV2 for open-vocabulary object detection, SAM for segmentation, and CLIP for semantic grounding.}
        \resumeItem{Deployed a pre-trained Vision-Language-Action (VLA) model for robot manipulation, executing policy inference from visual inputs and natural-language commands.}
        \resumeItem{Characterised generalisation and robustness failures of the VLA on unseen tasks, contributing concrete observations on where current VLA architectures break down.}
    \resumeItemListEnd

    
    \resumeProjectHeading{Multi Agent Path Finding in Unity3D $|$ \textnormal{\small DD2348 - KTH}}{Feb 2025 - May 2025}
    \resumeItemListStart 
        \resumeItem{Developed Hybrid A* path planner and waypoint tracking system in Unity3D for kinematically constrained single-agent (car/drone) navigation, achieving competitive performance.}
        \resumeItem{Engineered Conflict-Based Search (CBS) to solve MAPF for over 20 holonomic (drones) and non-holonomic (cars) agents, minimizing makespan in shared environments. }
        \resumeItem{Implemented a greedy dynamic goal assignment strategy for multi-agent teams ("bakery problem") integrated with CBS, optimizing task allocation and minimizing overall travel time. }
        % \resumeItem{battle snake }
    \resumeItemListEnd

    \newpage 
    \resumeProjectHeading{Reinforcement Learning for Pacman Capture The Flag $|$ \textnormal{\small DD2348 - KTH}}{Mar 2021 - Jun 2021}
    \resumeItemListStart 
        \resumeItem{Designed a hybrid AI for Pacman CTF (Unity3D), merging a custom Behavior Tree for high level strategy with RL models for specialized actions (attack, defend, escape).}
        \resumeItem{Trained RL (PPO) policies via Unity ML-Agents with role-specific designs (observations, rewards) and advanced methods (phased training, self-play, curriculum learning) for adaptive agent behavior.}
        \resumeItem{Evaluated strategic agent decision-making and adaptive behavior switching in a competitive, partially observable, dynamic multi-agent Pacman environment.}
    \resumeItemListEnd

    \resumeProjectHeading{Image analysis and Computer vision $|$ \textnormal{\small DD2348 - KTH}}{Oct 2024 - Nov 2024}
    \resumeItemListStart 
        \resumeItem{Applied Fourier Transforms for image filtering and frequency analysis; designed Gaussian smoothing filters via FFT to reduce noise and study effects on varying image resolutions.}
        \resumeItem{Implemented a multi-scale differential geometry edge detector (computing $L_{vv}, L_{vvv}$ for zero\-crossings)
        and used Hough Transform for line detection from edges.}
        \resumeItem{Implemented image matching (SIFT, RANSAC) to estimate homographies (planar scenes) and fundamental matrices, enabling 3D scene reconstruction via triangulation}  
    \resumeItemListEnd

    \resumeProjectHeading{Introduction to robotics $|$ \textnormal{\small DD2348 - KTH}}{Sep 2024 - Oct 2024}
    \resumeItemListStart 
        \resumeItem{Implemented an iterative inverse kinematics solver from scratch for a 7-DOF KUKA robot arm in ROS}
        \resumeItem{Applied RRT* to find collision-free paths for a Dubins car, showcasing path planning for non-holonomic vehicles.}
        \resumeItem{Engineered 2D occupancy grid mapping solutions in ROS, processing laser scans to build maps, identify free space, and generate C-spaces for navigation}  
        \resumeItem{Integrated a Behavior Tree mission planner for a TIAGo mobile manipulator (using ROS and Gazebo), enabling autonomous navigation, vision-based grasping, and robust task execution (e.g., kidnapped robot recovery).}
    \resumeItemListEnd
    
    % \resumeSubheading{Control of Automotive Systems}{Jun 2022 - Nov 2022}{\emph{Mentored by \href{https://home.iitm.ac.in/srikanthan/}{ Prof.Srikanthan Sridharan}}}{IIT Madras, India}
    \resumeProjectHeading{Control of Automotive Systems $|$ \textnormal{\small ED4060 - IIT Madras}}{Jun 2022 - Nov 2022}
    \resumeItemListStart
        \resumeItem{Designed a heading angle controller for autonomous ground vehicle considering its vehicle dynamics \& reduced error between the desired heading angle and the actual heading angle while working within given performance specifications}
        \resumeItem{Improved performance of pneumatic brake system by designing P and PI controllers to meet given specifications.}
        \resumeItem{Designed a nonlinear controller (Sliding Mode Control) to attenuate the effect of external forces on a tractor’s hydraulic hitch, to reduce excessive vibrations and to prevent front-wheel lift-off.}
    \resumeItemListEnd
    

\resumeSubHeadingListEnd


%-----------SKILLS-----------
\section{Technical Skills}

\begin{multicols}{2}
    \begin{itemize}[itemsep=- 1pt, parsep=3pt]
    \small{ 
              \item \textbf{Languages} : Python, C, C++, C\texttt{\#}\\ 
              \item \textbf{Robotics} : ROS2, Unity3D, MATLAB, MuJoCo\\ 
              \item \textbf{DL} : Pytorch, OpenCV, wandb, JAX \\ 
              \item \textbf{Tools} : VS Code, Docker, Git, \LaTeX, CUDA, Slurm
              \item \textbf{CAD} : Fusion360, Solidworks, Ansys, COMSOL
        } 
    \end{itemize}
\end{multicols}

%------RELEVANT COURSEWORK-------
\section{Relevant Coursework}

\textbf{Robotics} - Introduction to Field and Service Robotics, Mechanics of Serial Robots, \\
\textbf{Controls} - Modelling of Dynamical Systems, Control Theory Advanced \& Practice, Controls of Automotive Systems\\ 
\textbf{Programming} - Image Analysis and Computer Vision, Fundamentals of Deep Learning, Data Structures and Algorithms\\

\end{document}